\documentclass{amsart}
% For dealing with references we use the comment environment
\usepackage{verbatim}
\newenvironment{reference}{\comment}{\endcomment}
%\newenvironment{reference}{}{}
\newenvironment{slogan}{\comment}{\endcomment}
\newenvironment{history}{\comment}{\endcomment}

% For commutative diagrams we use Xy-pic
\usepackage[all]{xy}

% We use 2cell for 2-commutative diagrams.
\xyoption{2cell}
\UseAllTwocells

% We use multicol for the list of chapters between chapters
\usepackage{multicol}

% This is generally recommended for better output
\usepackage{lmodern}
\usepackage[T1]{fontenc}

% For cross-file-references

% Package for hypertext links:
\usepackage{hyperref}

% For any local file, say "hello.tex" you want to link to please

% Theorem environments.
%
\theoremstyle{plain}
\newtheorem{theorem}[subsection]{Theorem}
\newtheorem{proposition}[subsection]{Proposition}
\newtheorem{lemma}[subsection]{Lemma}

\theoremstyle{definition}
\newtheorem{definition}[subsection]{Definition}
\newtheorem{example}[subsection]{Example}
\newtheorem{exercise}[subsection]{Exercise}
\newtheorem{situation}[subsection]{Situation}

\theoremstyle{remark}
\newtheorem{remark}[subsection]{Remark}
\newtheorem{remarks}[subsection]{Remarks}

\numberwithin{equation}{subsection}

% Macros
%
\def\lim{\mathop{\mathrm{lim}}\nolimits}
\def\colim{\mathop{\mathrm{colim}}\nolimits}
\def\Spec{\mathop{\mathrm{Spec}}}
\def\Hom{\mathop{\mathrm{Hom}}\nolimits}
\def\Ext{\mathop{\mathrm{Ext}}\nolimits}
\def\SheafHom{\mathop{\mathcal{H}\!\mathit{om}}\nolimits}
\def\SheafExt{\mathop{\mathcal{E}\!\mathit{xt}}\nolimits}
\def\Sch{\mathit{Sch}}
\def\Mor{\mathop{\mathrm{Mor}}\nolimits}
\def\Ob{\mathop{\mathrm{Ob}}\nolimits}
\def\Sh{\mathop{\mathit{Sh}}\nolimits}
\def\NL{\mathop{N\!L}\nolimits}
\def\CH{\mathop{\mathrm{CH}}\nolimits}
\def\proetale{{pro\text{-}\acute{e}tale}}
\def\etale{{\acute{e}tale}}
\def\QCoh{\mathit{QCoh}}
\def\Ker{\mathop{\mathrm{Ker}}}
\def\Im{\mathop{\mathrm{Im}}}
\def\Coker{\mathop{\mathrm{Coker}}}
\def\Coim{\mathop{\mathrm{Coim}}}

% Boxtimes
%
\DeclareMathSymbol{\boxtimes}{\mathbin}{AMSa}{"02}

%
% Macros for moduli stacks/spaces
%
\def\QCohstack{\mathcal{QC}\!\mathit{oh}}
\def\Cohstack{\mathcal{C}\!\mathit{oh}}
\def\Spacesstack{\mathcal{S}\!\mathit{paces}}
\def\Quotfunctor{\mathrm{Quot}}
\def\Hilbfunctor{\mathrm{Hilb}}
\def\Curvesstack{\mathcal{C}\!\mathit{urves}}
\def\Polarizedstack{\mathcal{P}\!\mathit{olarized}}
\def\Complexesstack{\mathcal{C}\!\mathit{omplexes}}
% \Pic is the operator that assigns to X its picard group, usage \Pic(X)
% \Picardstack_{X/B} denotes the Picard stack of X over B
% \Picardfunctor_{X/B} denotes the Picard functor of X over B
\def\Pic{\mathop{\mathrm{Pic}}\nolimits}
\def\Picardstack{\mathcal{P}\!\mathit{ic}}
\def\Picardfunctor{\mathrm{Pic}}
\def\Deformationcategory{\mathcal{D}\!\mathit{ef}}

\setlength{\emergencystretch}{3em}
\begin{document}
\title[Stacks Project, Tag 0DEN]{656 --- K-flatness can be tested at points of a ringed site}
\maketitle
\noindent Copyright (C) 2005--2025 Johan de Jong. Stacks Project authors. Source: \href{https://stacks.math.columbia.edu/tag/0DEN}{Tag 0DEN}.

\noindent Editorial difficulty note (not author text). Difficulty H2: The difficult direction replaces a site by one with finite limits, spreads finitely presented stalk modules to neighborhoods, and transports acyclicity through K-flat resolutions. This unbounded derived-tensor argument with point-neighborhood approximation is beyond ordinary qualification homological algebra..

\noindent Editorial provenance note (not author text). History: excerpt from revision \texttt{a0444\allowbreak 6e57e\allowbreak c1fbc\allowbreak 252a8\allowbreak 71afc\allowbreak ec775\allowbreak 2fb28\allowbreak 07b14}, sites-cohomology.tex, lines 3766--3831. Reference presentation was adapted; original-author context and core support blocks were retained or added. The mathematical wording is unchanged. Licensed under GNU FDL 1.2 or later; no invariant sections or cover texts. See ../licenses/COPYING. Context blocks reproduce author definitions and supporting results; they are not additional counted problems.

\begin{definition}
\label{more-algebra-definition-K-flat}
Let $R$ be a ring. A complex $K^\bullet$ is called {\it K-flat}
if for every acyclic complex $M^\bullet$ the total complex
$\text{Tot}(M^\bullet \otimes_R K^\bullet)$ is acyclic.
\end{definition}

\begin{lemma}
\label{lemma-pullback-K-flat}
Let $f : (\Sh(\mathcal{C}'), \mathcal{O}') \to (\Sh(\mathcal{C}), \mathcal{O})$
be a morphism of ringed topoi. Let $\mathcal{K}^\bullet$ be a K-flat complex
of $\mathcal{O}$-modules whose terms are flat $\mathcal{O}$-modules. Then
$f^*\mathcal{K}^\bullet$ is a K-flat complex of $\mathcal{O}'$-modules whose
terms are flat $\mathcal{O}'$-modules.
\end{lemma}

\begin{proof}
The terms $f^*\mathcal{K}^n$ are flat $\mathcal{O}'$-modules by
Modules on Sites, Lemma \href{https://stacks.math.columbia.edu/tag/05VD}{lemma pullback flat (Tag 05VD)}.
Choose a diagram
$$
\xymatrix{
\mathcal{K}_1^\bullet \ar[d] \ar[r] &
\mathcal{K}_2^\bullet \ar[d] \ar[r] & \ldots \\
\tau_{\leq 1}\mathcal{K}^\bullet \ar[r] &
\tau_{\leq 2}\mathcal{K}^\bullet \ar[r] & \ldots
}
$$
as in Lemma \href{https://stacks.math.columbia.edu/tag/077J}{lemma resolution by direct sums extensions by zero (Tag 077J)}.
We will use all of the properties stated in the
lemma without further mention. Each $\mathcal{K}_n^\bullet$ is a bounded
above complex of flat modules, see
Modules on Sites, Lemma \href{https://stacks.math.columbia.edu/tag/03EV}{lemma j shriek flat (Tag 03EV)}.
Consider the short exact sequence of complexes
$$
0 \to \mathcal{M}^\bullet \to
\colim \mathcal{K}_n^\bullet \to
\mathcal{K}^\bullet \to 0
$$
defining $\mathcal{M}^\bullet$. By Lemmas \href{https://stacks.math.columbia.edu/tag/06YQ}{lemma bounded flat K flat (Tag 06YQ)} and
\href{https://stacks.math.columbia.edu/tag/06YR}{lemma colimit K flat (Tag 06YR)} the complex $\colim \mathcal{K}_n^\bullet$
is K-flat and by Modules on Sites, Lemma \href{https://stacks.math.columbia.edu/tag/03EU}{lemma colimits flat (Tag 03EU)}
it has flat terms. By Modules on Sites, Lemma \href{https://stacks.math.columbia.edu/tag/03EY}{lemma flat ses (Tag 03EY)}
$\mathcal{M}^\bullet$ has flat terms, by
Lemma \href{https://stacks.math.columbia.edu/tag/0G7B}{lemma K flat two out of three ses (Tag 0G7B)}
$\mathcal{M}^\bullet$ is K-flat, and by the long exact
cohomology sequence $\mathcal{M}^\bullet$ is acyclic (because
the second arrow is a quasi-isomorphism). The pullback
$f^*(\colim \mathcal{K}_n^\bullet) = \colim f^*\mathcal{K}_n^\bullet$
is a colimit of bounded below complexes of flat $\mathcal{O}'$-modules
and hence is K-flat (by the same lemmas as above).
The pullback of our short exact sequence
$$
0 \to f^*\mathcal{M}^\bullet \to
f^*(\colim \mathcal{K}_n^\bullet) \to
f^*\mathcal{K}^\bullet \to 0
$$
is a short exact sequence of complexes by
Modules on Sites, Lemma \href{https://stacks.math.columbia.edu/tag/0G6R}{lemma pullback ses (Tag 0G6R)}.
Hence by Lemma \href{https://stacks.math.columbia.edu/tag/0G7B}{lemma K flat two out of three ses (Tag 0G7B)}
it suffices to show that $f^*\mathcal{M}^\bullet$
is K-flat. This reduces us to the case discussed in the next paragraph.

\medskip\noindent
Assume $\mathcal{K}^\bullet$ is acyclic as well as K-flat and
with flat terms. Then Lemma \href{https://stacks.math.columbia.edu/tag/0G7C}{lemma K flat flat acyclic (Tag 0G7C)}
guarantees that all terms of $\tau_{\leq n}\mathcal{K}^\bullet$
are flat $\mathcal{O}$-modules. We choose a diagram as above and
we will use all the properties proven above for this diagram.
Denote $\mathcal{M}_n^\bullet$ the kernel of the map of complexes
$\mathcal{K}_n^\bullet \to \tau_{\leq n}\mathcal{K}^\bullet$
so that we have short exact sequences of complexes
$$
0 \to \mathcal{M}_n^\bullet \to \mathcal{K}_n^\bullet \to
\tau_{\leq n}\mathcal{K}^\bullet \to 0
$$
By Modules on Sites, Lemma \href{https://stacks.math.columbia.edu/tag/03EY}{lemma flat ses (Tag 03EY)}
we see that the terms of the complex $\mathcal{M}_n^\bullet$ are flat.
Hence we see that $\mathcal{M} = \colim \mathcal{M}_n^\bullet$
is a filtered colimit of bounded below complexes of flat modules
in this case. Thus $f^*\mathcal{M}^\bullet$ is K-flat
(same argument as above) and we win.
\end{proof}


\begin{lemma}
\label{lemma-check-K-flat-stalks}
Let $(\mathcal{C}, \mathcal{O})$ be a ringed site.
Let $\mathcal{K}^\bullet$ be a complex of $\mathcal{O}$-modules.
\begin{enumerate}
\item If $\mathcal{K}^\bullet$ is K-flat, then for every point $p$
of the site $\mathcal{C}$ the complex of $\mathcal{O}_p$-modules
$\mathcal{K}_p^\bullet$ is K-flat in the sense of
More on Algebra, Definition \ref{more-algebra-definition-K-flat}
\item If $\mathcal{C}$ has enough points, then the converse is true.
\end{enumerate}
\end{lemma}

\begin{proof}
Proof of (2). If $\mathcal{C}$ has enough points and
$\mathcal{K}_p^\bullet$ is K-flat for all points $p$ of $\mathcal{C}$
then we see that $\mathcal{K}^\bullet$ is K-flat because $\otimes$ and
direct sums commute with taking stalks and because we can check exactness
at stalks, see
Modules on Sites, Lemma \href{https://stacks.math.columbia.edu/tag/05V3}{lemma check exactness stalks (Tag 05V3)}.

\medskip\noindent
Proof of (1). Assume $\mathcal{K}^\bullet$ is K-flat.
Choose a quasi-isomorphism $a : \mathcal{L}^\bullet \to \mathcal{K}^\bullet$
such that $\mathcal{L}^\bullet$ is K-flat with flat terms, see
Lemma \href{https://stacks.math.columbia.edu/tag/06YS}{lemma K flat resolution (Tag 06YS)}. Any pullback
of $\mathcal{L}^\bullet$ is K-flat, see
Lemma \ref{lemma-pullback-K-flat}. In particular the stalk
$\mathcal{L}_p^\bullet$ is a K-flat complex of $\mathcal{O}_p$-modules.
Thus the cone $C(a)$ on $a$ is a K-flat
(Lemma \href{https://stacks.math.columbia.edu/tag/07A3}{lemma K flat two out of three (Tag 07A3)})
acyclic complex of $\mathcal{O}$-modules and it suffuces
to show the stalk of $C(a)$ is K-flat
(by More on Algebra, Lemma \href{https://stacks.math.columbia.edu/tag/06Y2}{lemma K flat two out of three (Tag 06Y2)}).
Thus we may assume that $\mathcal{K}^\bullet$ is K-flat and acyclic.

\medskip\noindent
Assume $\mathcal{K}^\bullet$ is acyclic and K-flat. Before continuing
we replace the site $\mathcal{C}$ by another one as in
Sites, Lemma \href{https://stacks.math.columbia.edu/tag/03CI}{lemma topos good site (Tag 03CI)}
to insure that $\mathcal{C}$ has all finite
limits. This implies the category of neighbourhoods of $p$ is filtered
(Sites, Lemma \href{https://stacks.math.columbia.edu/tag/00YB}{lemma neighbourhoods directed (Tag 00YB)})
and the colimit defining the stalk of a sheaf is filtered.
Let $M$ be a finitely presented $\mathcal{O}_p$-module.
It suffices to show that $\mathcal{K}^\bullet \otimes_{\mathcal{O}_p} M$
is acyclic, see
More on Algebra, Lemma \href{https://stacks.math.columbia.edu/tag/0E8F}{lemma universally acyclic K flat (Tag 0E8F)}.
Since $\mathcal{O}_p$ is the filtered colimit of $\mathcal{O}(U)$
where $U$ runs over the neighbourhoods of $p$, we
can find a neighbourhood $(U, x)$ of $p$ and a finitely
presented $\mathcal{O}(U)$-module $M'$ whose base change
to $\mathcal{O}_p$ is $M$, see
Algebra, Lemma \href{https://stacks.math.columbia.edu/tag/05N7}{lemma colimit category fp modules (Tag 05N7)}.
By Lemma \href{https://stacks.math.columbia.edu/tag/0E8K}{lemma restriction K flat (Tag 0E8K)}
we may replace $\mathcal{C}, \mathcal{O}, \mathcal{K}^\bullet$
by $\mathcal{C}/U, \mathcal{O}_U, \mathcal{K}^\bullet|_U$.
We conclude that we may assume there exists an $\mathcal{O}$-module
$\mathcal{F}$ such that $M \cong \mathcal{F}_p$.
Since $\mathcal{K}^\bullet$ is K-flat and acyclic,
we see that $\mathcal{K}^\bullet \otimes_\mathcal{O} \mathcal{F}$
is acyclic (as it computes the derived tensor product by
definition). Taking stalks is an exact functor, hence we get that
$\mathcal{K}^\bullet \otimes_{\mathcal{O}_p} M$
is acyclic as desired.
\end{proof}
\end{document}
